\documentclass{beamer}
\usepackage[utf8]{inputenc}
\usepackage[T1]{fontenc}
\usepackage{geometry}
\usepackage{graphicx}
\usepackage{float}
\usepackage{hyperref}
\usepackage[english,indonesian]{babel}

\title{Optimasi Pemrosesan Template LaTeX Menggunakan Pendekatan Berbasis AI}
\author{Academic Support Agent}
\date{\today}

\begin{document}
\maketitle

\begin{abstract}

\end{abstract}

\chapter{PENDAHULUAN}


\chapter{TINJAUAN PUSTAKA}


% Original content (commented out):
% \documentclass{beamer}

\title[Short title]{LSCM xyz:
big title}
\subtitle{Lecture 1 subtitle}
\date{\today}
\author{Dr. x\\
\texttt{mail@mm.com}}

\makeatletter
\def\input@path{{$beamerThemeHBKU$/}}
\makeatother
\usetheme{HBKU}


\begin{document}

\begin{frame}[plain]
\titlepage
\end{frame}

% frames with framesubitite
\begin{frame}[t]
\frametitle{There Is No Largest Prime Number}
\framesubtitle{Theorem and Proof}

\begin{theorem}[Example]
There is no largest prime number. 
\end{theorem}

\begin{proof}
The invertible elements in a field form a group under multiplication.
In particular, the elements
\begin{equation*}
1, 2, \ldots, p - 1 \in \mathbb{Z}_p
\end{equation*}
form a group under multiplication modulo~\(p\).
This is a group of order \(p - 1\).
For \(a \in \mathbb{Z}_p\) and \(a \neq 0\) we thus get \(a^{p-1} = 1 \in \mathbb{Z}_p\).
The claim follows.
\end{proof}
\end{frame}

\section{Section name}
\subsection{Subection name}

\begin{frame}[t]
\frametitle{The environments}
\framesubtitle{Lemma, Proposition, Corrollary and Definition}

\begin{lemma}[Title of lemma]
There is no largest prime number.
\end{lemma}

\begin{proposition}[Title of proposition]
There is no largest prime number.
\end{proposition}

\begin{corrollary}[Title of corrollary]
There is no largest prime number.
\end{corrollary}
\begin{definition}[Title of definition]
There is no largest prime number.
\end{definition}
\end{frame}


\begin{frame}{Enumerate environment}
\begin{enumerate}
 \item Suppose $p$ were the largest prime number.
 \item Let $q$ be the product of the first $p$ numbers.
 \item Then $q+1$ is not divisible by any of them.
 \item But $q + 1$ is greater than $1$, thus divisible by some prime
 number not in the first $p$ numbers.
 \end{enumerate}
\end{frame}



% frames without framesubitite
\begin{frame}[t]{Itemize environment}
 \begin{itemize}
\item one
\item two
\end{itemize}
\end{frame}

\begin{frame}{Figure, example and alert block}
\begin{figure}[htb]\centering
  \includegraphics[scale=0.25]{example-image-a}
\caption{Caption of figure}
\label{f11}
\end{figure}

\begin{example}
\begin{itemize}
\item
Lists change colour after the environment.
\end{itemize}
\end{example}

\begin{alertblock}{Important message}
        If a lot of text should be \alert{highlighted}, it is a good idea to put it in a box.
    \end{alertblock}
\end{frame}

\begin{frame}{Example block}
\begin{block}{Title of block}
If a lot of text should be highlighted, it is a good idea to put it in a box.
\end{block}
\end{frame}

\begin{frame}[fragile]{Example table}

{\renewcommand{\arraystretch}{1.5}
\arrayrulecolor{white}
\begin{table}[htb]\small
\caption{Table caption}
\label{t1}
\begin{tabular}{|>{\centering}m{0.3\textwidth}|
>{\raggedright}m{0.3\textwidth}|m{0.3\textwidth}<{\raggedleft}|}\hline
\rowcolor{HMblue} \textcolor{white}{Text center}&  
\textcolor{white}{Text left align}&  
\textcolor{white}{Text right align}\\\hline
\rowcolor{black!20}&  &  \\\hline
\rowcolor{black!10}&  &  \\\hline
\end{tabular}
\end{table}
}
\end{frame}

\end{document} 


\end{document}